\section{Topoi}
\label{section-topoi}

\noindent
Here is a definition of a topos which is suitable for our purposes.
Namely, a topos is the category of sheaves on a site. In order to specify
a topos you just specify the site. The real difference between a topos
and a site lies in the definition of morphisms. Namely, it turns out that
there are lots of morphisms of topoi which do not come from morphisms
of the underlying sites.

\begin{definition}[Topoi]
\label{definition-topos}
A {\it topos} is the category $\Sh(\mathcal{C})$ of sheaves
on a site $\mathcal{C}$.
\begin{enumerate}
\item Let $\mathcal{C}$, $\mathcal{D}$ be sites.
A {\it morphism of topoi} $f$ from $\Sh(\mathcal{D})$
to $\Sh(\mathcal{C})$ is given by a pair of functors
$f_* : \Sh(\mathcal{D}) \to \Sh(\mathcal{C})$
and
$f^{-1} : \Sh(\mathcal{C}) \to \Sh(\mathcal{D})$
such that
\begin{enumerate}
\item we have
$$
\Mor_{\Sh(\mathcal{D})}(f^{-1}\mathcal{G}, \mathcal{F})
=
\Mor_{\Sh(\mathcal{C})}(\mathcal{G}, f_*\mathcal{F})
$$
bifunctorially, and
\item the functor $f^{-1}$ commutes with finite limits, i.e.,
is left exact.
\end{enumerate}
\item Let $\mathcal{C}$, $\mathcal{D}$, $\mathcal{E}$ be sites.
Given morphisms of topoi
$f :\Sh(\mathcal{D}) \to \Sh(\mathcal{C})$ and
$g :\Sh(\mathcal{E}) \to \Sh(\mathcal{D})$ the
{\it composition $f\circ g$} is the morphism of topoi defined
by the functors
$(f \circ g)_* = f_* \circ g_*$ and
$(f \circ g)^{-1} = g^{-1} \circ f^{-1}$.
\end{enumerate}
\end{definition}

\noindent
Suppose that $\alpha : \mathcal{S}_1 \to \mathcal{S}_2$ is an equivalence of
(possibly ``big'') categories. If $\mathcal{S}_1$, $\mathcal{S}_2$
are topoi, then setting $f_* = \alpha$ and $f^{-1}$ equal to a quasi-inverse
of $\alpha$ gives a morphism $f : \mathcal{S}_1 \to \mathcal{S}_2$ of topoi.
Moreover this morphism is an equivalence in the $2$-category of topoi (see
Section \ref{section-2-category}).
Thus it makes sense to say ``$\mathcal{S}$ is a topos'' if $\mathcal{S}$
is equivalent to the category of sheaves on a site (and not necessarily
equal to the category of sheaves on a site). We will occasionally
use this abuse of notation.

\medskip\noindent
The {\it empty topos} is the topos
of sheaves on the site $\mathcal{C}$, where $\mathcal{C}$
is the empty category. We will sometimes write $\emptyset$ for this site.
This is a site which has a unique sheaf (since $\emptyset$ has no
objects). Thus $\Sh(\emptyset)$ is equivalent to the category
having a single object and a single morphism.

\medskip\noindent
The {\it punctual topos} is the topos of sheaves on the site
$\mathcal{C}$ which has a single object $pt$ and one morphism
$\text{id}_{pt}$ and whose only covering is the covering
$\{\text{id}_{pt}\}$. We will simply write $pt$ for this site.
It is clear that the category of sheaves $ = $ the category of
presheaves $ = $ the category of sets.
In a formula $\Sh(pt) = \textit{Sets}$.

\medskip\noindent
Let $\mathcal{C}$ and $\mathcal{D}$ be sites. Let
$f : \Sh(\mathcal{D}) \to \Sh(\mathcal{C})$ be a morphism of topoi.
Note that $f_*$ commutes with all limits and that
$f^{-1}$ commutes with all colimits, see Categories,
Lemma \ref{categories-lemma-adjoint-exact}.
In particular, the condition on $f^{-1}$ in the definition above
guarantees that $f^{-1}$ is exact. Morphisms of topoi are often constructed
using either Lemma \ref{lemma-cocontinuous-morphism-topoi}
or the following lemma.

\begin{lemma}
\label{lemma-morphism-sites-topoi}
Given a morphism of sites $f : \mathcal{D} \to \mathcal{C}$
corresponding to the functor $u : \mathcal{C} \to \mathcal{D}$
the pair of functors $(f^{-1} = u_s, f_* = u^s)$ is a morphism of topoi
$\Sh(\mathcal{D}) \to \Sh(\mathcal{C})$.
\end{lemma}

\begin{proof}
This is obvious from Definition \ref{definition-morphism-sites}.
\end{proof}

\begin{remark}
\label{remark-pt-topos}
There are many sites that give rise to the topos $\Sh(pt)$.
A useful example is the following. Suppose that $S$ is a set (of sets)
which contains at least one nonempty element. Let $\mathcal{S}$ be the
category whose objects are elements of $S$ and whose morphisms are
arbitrary set maps. Assume that $\mathcal{S}$ has fibre products.
For example this will be the case if $S = \mathcal{P}(\text{infinite set})$
is the power set of any infinite set (exercise in set theory).
Make $\mathcal{S}$ into a site by declaring
surjective families of maps to be coverings (and choose
a suitable sufficiently large set of covering families as in
Sets, Section \ref{sets-section-coverings-site}).
We claim that $\Sh(\mathcal{S})$ is equivalent to the category of
sets.

\medskip\noindent
We first prove this in case $S$ contains $e \in S$ which is a singleton.
In this case, there is an equivalence of topoi
$i : \Sh(pt) \to \Sh(\mathcal{S})$ given by
the functors
\begin{equation}
\label{equation-sheaves-pt-sets}
i^{-1}\mathcal{F} = \mathcal{F}(e), \quad
i_*E = (U \mapsto \Mor_{\textit{Sets}}(U, E))
\end{equation}
Namely, suppose that $\mathcal{F}$ is a sheaf on $\mathcal{S}$.
For any $U \in \Ob(\mathcal{S}) = S$ we can find a covering
$\{\varphi_u : e \to U\}_{u \in U}$, where $\varphi_u$
maps the unique element of $e$ to $u \in U$. The sheaf condition
implies in this case that
$\mathcal{F}(U) = \prod_{u \in U} \mathcal{F}(e)$.
In other words
$\mathcal{F}(U) = \Mor_{\textit{Sets}}(U, \mathcal{F}(e))$.
Moreover, this rule is compatible with restriction mappings. Hence
the functor
$$
i_* :
\textit{Sets} = \Sh(pt)
\longrightarrow
\Sh(\mathcal{S}), \quad
E \longmapsto (U \mapsto \Mor_{\textit{Sets}}(U, E))
$$
is an equivalence of categories, and its inverse is the functor
$i^{-1}$ given above.

\medskip\noindent
If $\mathcal{S}$ does not contain a singleton, then the functor
$i_*$ as defined above still makes sense. To show that it is still
an equivalence in this case, choose any nonempty $\tilde e \in S$
and a map $\varphi : \tilde e \to \tilde e$ whose image is a singleton.
For any sheaf $\mathcal{F}$ set
$$
\mathcal{F}(e) :=
\Im(
\mathcal{F}(\varphi) :
\mathcal{F}(\tilde e)
\longrightarrow
\mathcal{F}(\tilde e)
)
$$
and show that this is a quasi-inverse to $i_*$. Details omitted.
\end{remark}

\begin{remark}
\label{remark-morphism-topoi-big}
(Set theoretical issues related to morphisms of topoi. Skip
on a first reading.)
A morphism of topoi as defined above is not a set but a class.
In other words it is given by a mathematical formula rather
than a mathematical object. Although we may contemplate
the collection of all morphisms between two given topoi,
it is not a good idea to introduce it as a mathematical object.
On the other hand, suppose $\mathcal{C}$ and $\mathcal{D}$ are
given sites. Consider a functor
$\Phi : \mathcal{C} \to \Sh(\mathcal{D})$.
Such a thing is a set, in other words, it is a mathematical object.
We may, in succession, ask the following questions on $\Phi$.
\begin{enumerate}
\item Is it true, given a sheaf $\mathcal{F}$ on $\mathcal{D}$,
that the rule
$U \mapsto \Mor_{\Sh(\mathcal{D})}(\Phi(U), \mathcal{F})$
defines a sheaf on $\mathcal{C}$? If so, this defines a functor
$\Phi_* : \Sh(\mathcal{D}) \to \Sh(\mathcal{C})$.
\item Is it true that $\Phi_*$ has a left adjoint? If so,
write $\Phi^{-1}$ for this left adjoint.
\item Is it true that $\Phi^{-1}$ is exact?
\end{enumerate}
If the last question still has the answer ``yes'', then we obtain
a morphism of topoi $(\Phi_*, \Phi^{-1})$. Moreover, given any
morphism of topoi $(f_*, f^{-1})$ we may set
$\Phi(U) = f^{-1}(h_U^\#)$ and obtain a functor $\Phi$ as above
with $f_* \cong \Phi_*$ and $f^{-1} \cong \Phi^{-1}$ (compatible
with adjoint property).
The upshot is that by working with the collection of $\Phi$
instead of morphisms of topoi, we (a) replaced the notion of
a morphism of topoi by a mathematical object, and (b)
the collection of $\Phi$ forms a class (and not a collection
of classes). Of course, more can be said, for example one can work
out more precisely the significance of conditions (2) and (3) above;
we do this in the case of points of topoi in Section \ref{section-points}.
\end{remark}

\begin{remark}
\label{remark-quasi-continuous-morphism-topoi}
(Skip on first reading.)
Let $\mathcal{C}$ and $\mathcal{D}$ be sites.
A quasi-morphism of sites $f : \mathcal{D} \to \mathcal{C}$
(see Remark \ref{remark-quasi-continuous-morphism-sites})
gives rise to a morphism of topoi $f$ from
$\Sh(\mathcal{D})$ to $\Sh(\mathcal{C})$
exactly as in Lemma \ref{lemma-morphism-sites-topoi}.
\end{remark}













